% oct13-lem-separated-brackets: oct13.tex:1084-1091 (label lem:separated-brackets)
\paragraph{Paper excerpt.}
\begin{verbatim}
\begin{lemma}\label{lem:separated-brackets}
Let $\gauge$ be the Euclidean norm, and let $0<\drift<\nf1d$. There exist $0<c_0(d,\drift)\leq C_0(d,\drift)$ and $C(d,\drift)<\infty$ such that, for all sufficiently large~$n$, on~$\Omega_n$, for $1\leq i\leq m$ and $1\leq j\leq m$ with $j\neq i$,
\begin{equation*}
c_0\leq\langle S^i\rangle_n\leq C_0
\qquad\text{and}\qquad
|\langle S^i,S^j\rangle_n|\leq Cr_n^{-\nf14}\, .
\end{equation*}
\end{lemma}
\end{verbatim}
\paragraph{Lean statement.}
\begin{verbatim}
theorem CERW.LimitShapes.separated_brackets {d : ℕ} (hd : 2 ≤ d) :
    let ωd : ℝ := (volume (Metric.ball (0 : EuclideanSpace ℝ (Fin d)) 1)).toReal
    ∀ ε : ℝ, 0 < ε → ε < 1 / (d : ℝ) →
    let r : ℕ → ℝ := fun n => ((d + 1) * n / (2 * d * ε * ωd)) ^ ((1 : ℝ) / (d + 1))
    let e₁ : Site d := LatticeProb.unit ⟨0, by omega⟩
    let g : Site d → ℝ := CERW.latticeKernel d
    let σ : ℕ → ℝ := fun n =>
      if d = 2 then Real.sqrt (r n * Real.log (r n)) else Real.sqrt (r n)
    let k : ℕ → ℕ := fun n => ⌊r n ^ ((1 : ℝ) / 8)⌋₊
    let m : ℕ → ℕ := fun n => ⌊r n ^ ((1 : ℝ) / 8)⌋₊
    let y : ℕ → ℕ → Site d := fun n i => ((i * ⌈r n ^ ((1 : ℝ) / 4)⌉₊ : ℕ) : ℤ) • e₁
    let f : ℕ → ℕ → Site d → ℝ := fun n i z =>
      if d = 2 then g (z - (y n i + (k n : ℤ) • e₁)) - g (z - y n i) else g (z - y n i)
    ∀ K : CERW.LimitShapes.EventConstants,
      0 < K.cCo → 0 ≤ K.CCo → 0 ≤ K.CLoc → 0 ≤ K.CQuad → 0 ≤ K.CMart → 0 < K.CCross →
    ∃ c₀ C₀ C : ℝ, 0 < c₀ ∧ c₀ ≤ C₀ ∧ 0 < C ∧ ∃ n₀ : ℕ,
      ∀ ξ : Site d → EuclideanSpace ℝ (Fin d),
        (∀ x : Site d, x ≠ 0 →
          CERW.IsSubgradient (fun v : EuclideanSpace ℝ (Fin d) => ‖v‖) (CERW.toSpace x) (ξ x)) →
        ξ 0 = 0 → ∀ (X : ℕ → Site d) (n : ℕ), n₀ ≤ n →
        CERW.Support.Norm.ContactEvent.PathTyping d X →
        CERW.LimitShapes.EventEstimates (fun v : EuclideanSpace ℝ (Fin d) => ‖v‖) ε ξ K X n →
        ∀ i j : ℕ, 1 ≤ i → i ≤ m n → 1 ≤ j → j ≤ m n →
            let B : ℝ := CERW.dynkinBracket (CERW.stepProb d ε) (f n i) (f n j) X n /
              σ n ^ 2
            (i = j → c₀ ≤ B ∧ B ≤ C₀) ∧ (i ≠ j → |B| ≤ C * r n ^ (-(1 : ℝ) / 4))
\end{verbatim}
\paragraph{Proof of the registered statement.}
Fix the event constants in the registered type and specialize the norm to
the Euclidean norm. On a legal path satisfying the event, the coarse upper
radius bound places the whole path strictly inside the stopping radius.
Thus the stopped quadratic process equals the quadratic martingale.
The coarse lower radius bound converts its error estimate to the quadratic
input with constant $K.\mathrm{CQuad}/K.\mathrm{cCo}$.
The contact implication (node \texttt{oct13-lem-contact}), the coarse bounds
and the reduced local approximation give the deterministic inner estimate
through \texttt{InnerApproximation.lean}.

The actual-bracket local martingale estimate gives the weighted pointwise
estimate as in the contact proof. The distance crossing constituent gives
the Euclidean crossing bound: for each level $b\geq0$, apply the inequality
to a point of maximum radius and use that its radius minus $b$ is at most
its distance to the closed ball. These inputs, the inner estimate, and the existing
outer-radius and local-profile arguments give a uniform profile error
$C\sqrt{r_n}(\log n)^{3/2}$; the higher-dimensional bound is stronger.
Apply the proved deterministic bracket calculation
\texttt{brackets\_core} in \texttt{CERW/Support/Lower/SeparatedBrackets.lean} to this profile.
Its specified translated kernels, spacing, normalization $\sigma_n^2$,
diagonal bounds and off-diagonal error are exactly those in the registered
conclusion. All time thresholds and constants are chosen before the field
and path, so taking their maximum gives the required uniform threshold.

This proof establishes the deterministic implication for fixed constants.
The proved node \texttt{lem-separated-brackets-common-event} supplies the
Euclidean common-event constants at $p=10$ and combines their probability
bound with this conclusion, preserving the source constant dependence.
